\chapter*{\thechapter \quad Revisão da Literatura}
\addcontentsline{toc}{chapter}{\thechapter \quad Revisão da Literatura}
\paragraph{}
Este capítulo apresenta os conceitos teóricos necessários ao desenvolvimento deste projeto, assim como a respetiva revisão da literatura.

A revisão da literatura foi realizada com o recurso a diversos serviços que auxiliaram na pesquisa de documentos relativos às áreas de estudo abordadas, sendo eles: \textit{Google Scholar}, \textit{Scopus}, \textit{IEEE} e \textit{Google}. O processo de seleção da literatura baseou-se sobretudo no tipo e relevância para o estudo que o documento possuía. Os tipos dos documentos selecionados surgem maioritariamente em formato de livro, artigo ou \textit{white paper}. A qualidade e relevância do documento é avaliada primeiramente pela análise do resumo do documento e, posteriormente, caso seja considerado oportuno para o estudo, do seu conteúdo.  

A revisão da literatura começou através da pesquisa pela palavra “\textit{Cibersegurança}", sendo que, os resultados obtidos remetiam para a existência de inúmeras perspetivas 